\documentclass{article}

% Symbols
\usepackage{recycle}
\usepackage{amsfonts, amsthm}
\usepackage{upgreek}
\usepackage{physics}
\usepackage{cancel}
\usepackage{amssymb, latexsym, amsmath}

% Proof
\renewcommand*{\proofname}{\textbf{Demostraci\'on:}}

% Graphics
\usepackage{graphicx}
\usepackage{pgf}

% Color a letras.
%
%\usepackage[usenames,dvipsnames,svgnames,table]{xcolor}

% Tikz
\usepackage{tikz}
\usetikzlibrary{arrows,automata}
\usepackage{tikz}
\usetikzlibrary{arrows,automata}

\usetikzlibrary{shapes,calc}
\tikzstyle{edge}=[shorten <=2pt, shorten >=2pt,
  >=stealth, line width=1.1pt]
\tikzstyle{blueE}=[shorten <=2pt, shorten >=2pt,
  >=stealth, line width=1.5pt, blue]
\tikzstyle{blackV}=[circle, fill=black,
  minimum size=6pt,
  inner sep=0pt, outer sep=0pt]
\tikzstyle{blueV}=[circle, fill=blue, draw,
  minimum size=6pt, line width=0.75pt,
  inner sep=0pt, outer sep=0pt]
\tikzstyle{redV}=[circle, fill=red, draw,
  minimum size=6pt, line width=0.75pt,
  inner sep=0pt, outer sep=0pt]
\tikzstyle{redSV}=[semicircle, fill=red, minimum
  size=3pt, inner sep=0pt, outer sep=0pt,
  rotate=225]
\tikzstyle{blueSV}=[semicircle, fill=blue, minimum
  size=3pt, inner sep=0pt, outer sep=0pt,
  rotate=225]
\tikzstyle{blackSV}=[semicircle, fill=black, minimum
  size=3pt, inner sep=0pt, outer sep=0pt,
  rotate=225]
\tikzstyle{vertex}=[circle, draw, minimum size=6pt,
  line width=0.75pt, inner sep=0pt,
  outer sep=0pt]

% Margins
\addtolength{\voffset}{-1cm}
\addtolength{\hoffset}{-1.5cm}
\addtolength{\textwidth}{3cm}
\addtolength{\textheight}{2cm}

%Header-Footer
\usepackage{fancyhdr}
\renewcommand{\headrulewidth}{1pt}

\newcommand{\set}[1]{%
  \left\{ #1 \right\}%
}

\pagenumbering{gobble}
\footskip = 50pt
\renewcommand{\headrulewidth}{1pt}

\pagestyle{fancyplain}

\begin{document}
\title{UNIVERSIDAD AUT\'ONOMA DE M\'EXICO\\ Facultad de Ciencias}
\author{Autores:
  \\ Fernanda Villaf\'an Flores
  \\ Fernando Alvarado Palacios
  \\ Adri\'an Aguilera Moreno}
\date{}
\maketitle
\begin{center}
  \includegraphics[scale=0.20]{../Imagen/Portada.jpg}\\[0.4cm]
  \Large
  \bf{Gr\'aficas y Juegos}
  \normalsize
\end{center}
\newpage
\fancyhead[r]{ Gr\'aficas y Juegos 2022-1}
\section*{\LARGE{Reposición}}

$$\text{\textbf{\large Ejercicio de la Tarea 2}}$$
\begin{enumerate}
  %%%%%%%%%%%%%%%%%%%%%%%%%%%%%%%%%%%%%%%%%%%%%%%%%% Ejercicio 3 de la Tarea 02.
  \item \text{[Ejercicio 3]}

    Sea $G$ una gr\'afica conexa. Demuestre que si $G$ no es completa, entonces contiente a $P_3$ como subgr\'afica inducida.

    \renewcommand\qedsymbol{QED}
    \begin{proof}

      Para este ejercicio necesitamos que $|V_G| \geq 3$, para las
      gr\'aficas que no cumplan esto se tendr\'a la demostraci\'on por vacuidad.
      Nótese que el hecho de que $G$ no sea completa implica que para al
      menos $x,y \in V_G$ se tiene que $xy \notin E_G$.

      \begin{center}
        \fbox{
          \begin{minipage}[b][1\height]%
            [t]{0.867\textwidth}
            Previo a la demostraci\'on, provemos que en una gr\'afica conexa
            siempre podemos construir una trayectoria con exactamente $3$
            v\'ertices:
            \vspace*{0.3cm}

            Sea $x \in V_G$, por definici\'on de conexidad y como $|V_G| \geq 3$,
            tenemos ha $x,y \in V_G$ tales que $xy \in E_G$, luego $x$ es vecino
            a alg\'un v\'ertice distinto a $y$ (o $y$ es vecino de alg\'un v\'ertice
            distinto de $x$), pues en caso contrario $xy$ ser\'ia una componente
            conexa contenida en $G$ y $xy \not= G$!! lo que contradice la hip\'otesis
            de que $G$ es conexa. Supongamos, sin p\'erdida de generalidad, que $z$
            es vecino de $x$ y $z \not= y$, luego $zxy$ es una trayectoria de orden
            exactamente $3$.

            \hfill $\square$
        \end{minipage}}
      \end{center}

      Para este ejercicio basta analizar $2$ posibles casos\footnote{Se
      analizan los casos ``extremos'', pues los casos intermedios son combinaciones
      de estos.}:

      \textcolor{blue}{Caso 1:} Si $G + e$ es completa, donde $e = xy-$arista para
      $x,y \in V_G$. Por \textbf{\textit{Prop.}} $\mathbf{1.64}$ y por hipótesis
      sabemos que existe un $xy$-camino en $G$, luego por \textbf{\textit{Prop.}}
      $\mathbf{1.62}$ sabemos que hay, en particular, una $xy-$trayectoria en $G$,
      luego hay alguna $xy-$trayectoria de orden $3$ (esto lo sabemos gracias al
      resultado mostrado previamente) y supongamos, sin p\'erdida de generalidad,
      que \'esta es $T = (x,z,y)$, para $z \in V_G$, notemos que $T$ tiene tamaño
      igual a $2$, pues existen las aristas $zx, zy$ pero no $xy$ (por como definimos
      este caso), luego $T$ es $P_3$ y concluimos que $P_3$ es subgr\'afica inducida
      de $G$.

      \textcolor{blue}{Caso 2:} Si $G$ es un \'arbol, esto nos indica que $G$ es
      $1-$conexa, y es por eso que se considera este caso como el m\'inimo para el
      que se cumplir\'a la condici\'on a demostrar. Sabemos por el teorema de
      caracterizaci\'on de \'arboles que cada arista en $G$ ser\'a un puente, y
      por el resultado previamente mostrado sabemos que existe una trayectoria $T$
      en $G$ de orden exactamente $3$, as\'i $T$ es claramente $P_3$ y concluimos
      $P_3$ es subg\'afica inducida de $G$.

      De los casos anteriores concluimos que el enunciado es verdadero.
    \end{proof}

    %%%%%%%%%%%%%%%%%%%%%%%%%%%%%%%%%%%%%%%%%%%%%%%%%% Ejercicio 1 extra de la Tarea 02.
    \item \text{[Ejercicio 1 extra]}

        Sea $G$ una gr\'afica. Demuestre que $G$ es $k$-partita completa si y s\'olo si no contiene a $K_{k+1}$ ni a $\overline{P_3}$ como subgr\'aficas inducidas.

        \begin{proof}
        En este ejercicio analizaremos 2 casos posibles:

            \begin{itemize}
                \item[$\Rightarrow$)] Procedamos reducción al absurdo .

                    \begin{itemize}
                        \item[$\cdot$)] Supongamos que $\overline{P_3}$ es subgr\'afica inducida
                        de $G$, por definición de $k$-partita completa $\overline{P_3}$ no está
                        en la misma parte (pues, en caso de estarlo hay una adyacencia en $2$
                        vértices de la misma parte). Luego $\overline{P_3}$ está en $2$ o $3$
                        partes distintas y habrá un $x \in \overline{P_3}$ que no se relacionará
                        con al menos $1$ vértice en algunas de las partes y por tanto, $G$ no es
                        $k$-partita completa!! (lo que no cumple es ser completa bajo el supuesto
                        tomado) y he aquí una contradicción de suponer que $\overline{P_3}$ es
                        subgr\'afica inducida de $G$. Por lo tanto, concluimos que $\overline{P_3}$
                        no es subgr\'afica inducida de $G$.

                        \item[$\cdot$)] Supongamos que $K_{k + 1}$ es subgr\'afica inducida de $G$,
                        entonces hay $1$ vértice de $K_{k + 1}$ en cada una de las partes (lo que
                        suma $k$ v\'ertices) y un $x \in K_{k + 1}$ en alguna parte tal que se
                        relaciona con exactamente un v\'ertice en esa parte y por tanto, $G$ no
                        es $k$-partita completa!! (no cumple el ser $k$-partita) y he aquí una
                        contradicción de suponer que $K_{k + 1}$ es subgr\'afica inducida de $G$.
                        Por lo tanto, concluimos que $K_{k + 1}$ no es subgr\'afica inducida de $G$.
                    \end{itemize}

                \item[$\Leftarrow$)]
            \end{itemize}
            De los casos anterior concluimos que $G$ es $k$-partita completa si y
            s\'olo si no contiene a $K_{k+1}$ ni a $\overline{P_3}$ como subgr\'aficas
            inducidas.
        \end{proof}

    $$\text{\textbf{\large Ejercicio de la Tarea 4}}$$
    %%%%%%%%%%%%%%%%%%%%%%%%%%%%%%%%%%%%%%%%%%%%%%%%%% Ejercicio 3 de la tarea 4

   \item \text{[Ejercicio extra 3]}

        Sea $\mathcal{T}$ una familia de sub\'arboles
        de un \'arbol $T$.   Deduzca, por inducci\'on sobre
        $|\mathcal{T}|$, que si cualesquiera dos elementos
        de $\mathcal{T}$ tienen un v\'ertice en com\'un,
        entonces hay un v\'ertice de $T$ que est\'a en
        todos los elementos de $\mathcal{T}$.

        \begin{proof}
            Demostración por induccion sobre el numero de subarboles

            Paso base: $\mathcal{T} = 3$

            Sean $T_1$, $T_2$, $T_3$ subarboles de T

            Pd Existe  un vértice x que pertenece a T tal que $T_1 \cap T_2=x \rightarrow  T_1 \cap T_2 \cap T_3 = x$

            Dem (Reduccion al absurdo)

            Supongamos que $T_1 \cap T_2 \cap T_3 = \emptyset$ (algo que no puede pasar es que $T_1 \cap T_2 = T_1 \cap T_3 = T_2 \cap T_3  = \emptyset$ por hipotesis T es conexa y tambien por hipotesis $T_1 \cap T_2$ diferente del vacio) $\rightarrow$ existe $x_1$ tal que $x_1$ pertenece a $T_1 \cap T_2$ y $x_1$ no pertenece a $T_3$, existe $x_2$ tal que $x_2$ pertenece a $T_1 \cap T_3$ y $x_2$ no pertenece a $T_2$, existe $x_3$ tal que $x_3$ pertenece a $T_2 \cap T_3$ y $x_3$ no pertenece a $T_1  \rightarrow x_1, x_2, x_3$ son vértices diferentes $\rightarrow$ por las intersecciones  y gracias que T  es conexo  $\rightarrow$ s.p.g. podemos formar un ciclo $(x_1, ... ,x_2, ... ,x_3, ...,x_1)$ (Lo que es una contradicción, ya que los arboles son aciclicos)

            Por lo tanto $T_1 \cap T_2 \cap T_3   \neq
            \emptyset \rightarrow$ existe un vértice compartido para los 3 subárboles

            Hipótesis de inducción: supongamos para $\mathcal{T} = k$

            Supongamos que si $\mathcal{T} = k$ tal que existen i,j tal que $i \neq j$ e i,j pertenezen a $\left\lbrace 1,...,k \right\rbrace$ donde $T_i \cap T_j \neq \emptyset \rightarrow  \bigcap_{r=1}^{k}{T_r} \neq \emptyset$

            Paso inductivo: Pd para $\mathcal{T} = k+1$

            Pd si existen i,j tal que $i \neq j$ e i,j pertenezen a $\left\lbrace 1,...,k+1 \right\rbrace$ donde $T_i \cap T_j \neq \emptyset \rightarrow  \bigcap_{r=1}^{k+1}{T_r} \neq \emptyset$

            Dem (Reducción al absurdo)

            Supongamos $\bigcap_{r=1}^{k+1}{T_r} = \emptyset$ y existen i,j pertenezen a $\left\lbrace 1,...,k+1 \right\rbrace$ donde $T_i \cap T_j \neq \emptyset \rightarrow $ consideremos a T' como el subárbol formado por la unión de todos los subárboles de $T_1, T_2,...;T_{k+1}$ menos los subárboles $T_i$ y $T_j$. Es decir T'= $\bigcup_{r=1, r \neq i,  r \neq j }^{k+1}{T_r} \rightarrow T' \cap T_i \cap T_j = \emptyset$, pero por paso base esto es una contradicción $\rightarrow$ existe un x tal que x pertenece a  $\bigcap_{r=1}^{k+1}{T_r}$

            Por lo tanto la porposición es verdadera.
        \end{proof}

  %%%%%%%%%%%%%%%%%%%%%%%%%%%%%%%%%%%%%%%%%%%%%%%%%% Tarea 5
  $$\text{\textbf{\large Ejercicio de la Tarea 5}}$$

  %%%%%%%%%%%%%%%%%%%%%%%%%%%%%%%%%%%%%%%%%%%%%%%%%% Preguta 1
  \item \text{[Pregunta 1]}

    Demuestre que si $G$ es simple y $3$-regular, entonces $\kappa = \kappa'$.

    \begin{proof}
      Sea G 3-regular $\rightarrow d(v)=3$ para todo v que pertenece a V $\rightarrow$ sea $v'$ un vertice de G  $\rightarrow$ para desconectar a $v'$ de G, solo basta  con "cortar" las 3 aristas de $v' \rightarrow$ k'=3.

      Por lo tanto $k=k'$.
    \end{proof}

  %%%%%%%%%%%%%%%%%%%%%%%%%%%%%%%%%%%%%%%%%%%%%%%%%% Pregunta 2
  \item \text{[Pregunta 2]}

    Demuestre que una gr\'afica es $2$-conexa por aristas si y s\'olo si
    cualesquiera dos v\'ertices est\'an conectados por al menos dos trayectorias
    ajenas por aristas.

    \begin{proof}
      $\Longrightarrow$) Sea G una gráfica 2-conexa $\rightarrow$ por teorema visto en clase en G existe  un cliclo C que contendra 2 vertices v y u donde u y v pertenecen a G $\rightarrow$ podemos tener la trayectoria $P=(v,C,u)$ pero como C es un ciclo $\rightarrow$ tambien existira la trayectoria $P'=(u,V,v)$

      Por lo tanto existen 2 trayectorias ajenas por aristas en una gráfica 2-conexa


      $\Longleftarrow$) Sean u y v cualesquiera vértices de una gráfica G y si u y v están conectados por dos trayectorias ajenas por aristas P y P' $\rightarrow$ si unimos  uPv y vP'v obtendremos un ciclo C que ira de uPvP'u $\rightarrow$ sea G' una gráfica igual al ciclo C, si borramos una arista a G' $\rightarrow$ G' seguira siendo conexa  $\rightarrow$ G' es 2-conexa $\rightarrow$ G será 2-conexa ya que para todo v y u que pertenecen G existen 2 trayectorias ajenas por vertices
    \end{proof}

  $$\text{\textbf{\large Ejercicio de la Tarea 7}}$$
  %%%%%%%%%%%%%%%%%%%%%%%%%%%%%%%%%%%%%%%%%%%%%%%%%% Ejercicio 1 de la tarea 7
  \item \text{[Ejercicio 1]}

    \begin{enumerate}
      %~~~~~~~~~~~~~~~~~~~~~~~~~~~~~~~~ inciso (a)
      \item Demuestre que si $G$ tiene di\'ametro mayor que $3$ (posiblemente
      infinito), entonces $\overline{G}$ tiene di\'ametro menor que $3$.
      Concluya que si $G$ es inconexa, entonces $\overline{G}$ es conexa.

        \begin{proof}
          Sea $G$ una gráfica.

          Probaremos que $\overline{G}$ tiene diámetro menor a $3$.

          Primero, sabemos que $G$ tiene diametro mayor a $3$ entonces tomemos una trayectoria $P$ de $G$ tal que su longitud es $n$ (con $n > 3$).

          La denotaremos como:
          $$P = (x_{0}, x_{1}, \dots, x_{n})$$

          Ahora, por definición de $\overline{G}$ es tal que:
          $$|E_{\overline{G}}| = \binom{|V| - 1}{2} - |E_{G}|$$

          Por tanto, en $\overline{G}$ la trayectoria $P$ cambia de la siguiente manera:
          \begin{itemize}
            \item El vértice $x_{0}$ es adyacente a los vértices $x_{2}, x_{3}, \dots, x_{n}$, donde esto equivale a $n - 1$ vértices.

            \item El vértice $x_{1}$ es adyacente a los vértices $x_{3}, x_{4}, \dots, x_{n}$, donde esto equivale a $n - 2$ vértices.

            \item El vértice $x_{2}$ es adyacente a los vértices $x_{0}, x_{4}, \dots, x_{n}$, donde esto equivale a $n - 2$ vértices.
          \end{itemize}

          Siguiendo este procedimiento, tenemos lo siguiente:
          \begin{itemize}
            \item El vértice $x_{i}$ es adyacente a los vértices $x_{i-2}, x_{i+2}, x_{i+3}, \dots, x_{n}$, con $i > 1$.
          \end{itemize}

          Así, notemos lo siguiente:

          En $\overline{G}$ $x_{0}$ \textbf{no es} adyacente a $x_{1}$, entonces necesitamos otro vértice $x_{3}$ para llegar a $x_{1}$. Lo que implica que nos toma distancia igual a $2$ llegar de $x_{0}$ a $x_{1}$.

          De la misma forma, $x_{1}$ \textbf{no es} adyacente a $x_{0}$ ni a $x_{2}$, entonces necesitamos otro vértice $x_{4}$ para llegar a $x_{0}$ o $x_{2}$. Lo que implica que nos toma distancia igual a $2$ llegar de $x_{1}$ a $x_{0}$ o de $x_{1}$ a $x_{2}$.

          De lo anterior obtenemos que:

          El vértice $x_{i}$ \textbf{no es} adyacente al vértice $x_{i-1}$ ni al vértice $x_{i+1}$, entonces necesitamos otro vértice $x_{i+2}$ para llegar a $x_{i-1}$ o $x_{i+1}$. Lo que implica que nos toma distancia igual a $2$ llegar de $x_{i}$ a $x_{i-1}$ o de $x_{i}$ a $x_{i+1}$.

          Por lo tanto, $\overline{G}$ tiene diámetro menor a $3$. \\
          \\
          \textbf{Por último}, probemos que si $G$ es inconexa entonces $\overline{G}$ es conexa.

          \begin{proof}
            Sean $u, v \in G$ cualesquiera.

            Supongamos que $G$ es inconexa.

            Sabemos que si $uv \notin E_{G}$, entonces $uv \in E_{\overline{G}}$.

            Por lo antes visto, tenemos que:

            Si $u$ es adyacente a $v$, como el diámetro de $\overline{G}$ es menor a $3$, entonecs existe otro vértice $w \in G$ tal que $uv, vw \in E_{\overline{G}}$.

            Es decir, podemos llegar desde un vértice $u$ a cualquier otro vértice $w$ usando al vértice $v$.

            Por lo tanto, ya que $u, v$ fueron arbitrarios, podemos concluir que $\overline{G}$ es conexa.
          \end{proof}

          Así, concluímos que si $G$ es inconexa, entonces $\overline{G}$ es conexa.
        \end{proof}
      %~~~~~~~~~~~~~~~~~~~~~~~~~~~~~~~~ inciso (b)
      \item Una gr\'afica $G$ es autocomplementaria si $G \cong \overline{G}$.
      Demuestre que si $G$ es autocomplementaria, entonces $|V|
      \stackrel{4}{\equiv} 0$ o $|V| \stackrel{4}{\equiv} 1$.

      \begin{proof}
        Notemos lo siguiente:
        $$|V| \stackrel{4}{\equiv} 1 \longrightarrow |V| \equiv 1 \mod{4} \text{\hspace{0.3cm} y \hspace{0.3cm}} |V| \stackrel{4}{\equiv} 0 \longrightarrow |V| \equiv 0 \mod{4}$$
        Recordando la definición de $\mod$, tenemos:
        \begin{align*}
          |V| \equiv 1 \mod{4}
          & \longrightarrow 4 \text{ \Big \vert } |V| - 1 \\
          & \longrightarrow |V| - 1 = 4 \cdot k, \text{ con $k \in \mathbb{N}$} \\
          & \longrightarrow |V| = 4 \cdot k + 1
        \end{align*}
        De la misma forma, tenemos:
        \begin{align*}
          |V| \equiv 0 \mod{4}
          & \longrightarrow 4 \text{ \Big \vert } |V| - 0 \\
          & \longrightarrow |V| = 4 \cdot k, \text{ con $k \in \mathbb{N}$} \\
          & \longrightarrow |V| = 4 \cdot k
        \end{align*}

        Sea $G$ una gráfica cualquiera.

        Como $G \cong \overline{G}$, entonces $|V_{G}| = |V_{\overline{G}}|$ y $|E_{G}| = |E_{\overline{G}}|$.

        Sea $n = |V_{G}|$. Tenemos que:
        \begin{center}
          $|E_{G}| = |E_{\overline{G}}| = \dfrac{|E(K_{n})|}{2}$

          , donde $K_{n}$ es la gráfica completa de $n$ vértices.
        \end{center}

        Ahora, recordemos que:
        \begin{center}
            \textit{Una gráfica es completa si su conjunto de aristas es igual a} $\displaystyle \binom{V}{2}$.
        \end{center}

        Así,
        \begin{align*}
          \binom{V}{2}
          & = \binom{n}{2} \\
          & = \dfrac{n!}{2! \cdot (n -2)!} \\
          & = \dfrac{n \cdot (n -1) \cdot (n - 2)!}{2 \cdot (n - 2)!} \\
          & = \dfrac{n(n - 1)}{2}
        \end{align*}

        Como $|E_{G}| = |E_{\overline{G}}|$, entonces sabemos que:
        \begin{align*}
          \dfrac{|E(K_{n})|}{2}
          & = \dfrac{1}{2} (|E(K_{n})|) \\
          & = \dfrac{1}{2} \left[\dfrac{n(n - 1)}{2} \right] \\
          & = \dfrac{n(n - 1)}{4}
        \end{align*}

        Por tanto, como $n$ y $(n - 1)$ son número enteros consecutivos, el máximo común divisor de ambos números es $1$.

        Luego, $4$ es divisor de $n$ o de $(n -1)$. Es decir, $n = 4k$ o $n - 1 = 4k \longrightarrow n = 4k + 1$ (con $k \in \mathbb{N}$).

        Por lo tanto, podemos concluir que si $G$ es autocomplementaria, entonces $|V| \stackrel{4}{\equiv} 0$ o $|V| \stackrel{4}{\equiv} 1$. \\
      \end{proof}
    \end{enumerate}
\end{enumerate}
\end{document}
